\documentclass[11pt,a4paper]{article}

\usepackage[margin=1in]{geometry}
\usepackage[T1]{fontenc}
\usepackage[utf8]{inputenc}
\usepackage{lmodern}
\usepackage{amsmath,amssymb}
\usepackage{graphicx}
\usepackage{booktabs,tabularx}
\usepackage{xcolor}
\usepackage[hidelinks]{hyperref}
\graphicspath{{figures/}}

% Visible drafting prompts: replace every occurrence before submission.
\newcommand{\draftnote}[1]{\par\smallskip\noindent\textcolor{blue}{\textbf{To write:} #1}\par\smallskip}

\title{CardioFusion: An Experimental Study of ECG--PCG Fusion\\for HFrEF Detection}
\author{Marwan El-Gharbawy \and Nadine Mohamed \and Shams Quraytum}
\date{\today}

\begin{document}
\maketitle

% \begin{abstract}
% abstract was removed as it's not an actual research paper
% abstract should be written last
% % ------------- abstract -------------
% \draftnote{Summarize the problem, dataset, compared models, main numerical
% finding, and principal limitation in one paragraph. Do not report unverified
% results or claim clinical readiness.}
% \end{abstract}
%
% ------------- introduction -------------
\section{Introduction}
\label{sec:introduction}

This project investigates whether the combination of electrocardiogram (ECG) and
phonocardiogram (PCG) recordings improves the detection of heart
failure with reduced ejection fraction (HFrEF). We compare unimodal models
with multimodal fusion approaches on a public paired-signal dataset.

ECG records the electrical activity of the heart, while PCG records heart
sounds associated with its mechanical activity. These signals may provide
complementary information for HFrEF screening, as investigated by
Ma et al.~\cite{ma2026hierarchical}. In this course project, we study this
possibility through experiments with individual modalities and their fusion.

% \draftnote{Introduce the clinical motivation with appropriate references.
% Explain what each modality measures and why complementary information might
% help classification. State the project's scope as an empirical course study.}

\subsection{Related work and study objective}
Ma et al.~\cite{ma2026hierarchical} introduced LGC$^{2}$-Net, a model that
combines ECG and PCG features at several stages and learns relationships
between the four recording sites. Their full model achieved a reported
five-fold AUROC of $0.960\pm0.020$, with sensitivity of
$(90.00\pm6.97)\%$ and specificity of $(92.73\pm2.62)\%$.

Using their public dataset, our study investigates whether combining ECG
and PCG improves detection compared with either modality alone. We compare
ECG-only and PCG-only models with concatenation and attention fusion.
Section~\ref{sec:comparison} presents our results alongside the published
values and explains the differences in evaluation.
% ------------- dataset preprocessing -------------
\section{Dataset and Preprocessing}
\label{sec:data}

\subsection{Dataset and prediction target}
The public dataset contains 620 patients and 2,480 paired ECG--PCG
recordings, with four recording sites per patient~\cite{hfrefdataset}.
Table~\ref{tab:cohorts} summarizes the development and community test
cohorts. A patient is labeled positive for HFrEF when left ventricular
ejection fraction (LVEF) is at or below 40\%. Patients with LVEF above
40\% form the negative class; this label does not necessarily mean that
they have no other cardiac condition.

\begin{table}[htbp]
\centering
\caption{Patient counts and HFrEF prevalence in the two cohorts.}
\label{tab:cohorts}
\begin{tabular}{lrrrr}
\toprule
Cohort & Patients & HFrEF & Non-HFrEF & HFrEF (\%) \\
\midrule
Development & 500 & 60 & 440 & 12.00 \\
Community test & 120 & 5 & 115 & 4.17 \\
\bottomrule
\end{tabular}
\end{table}

The four sites are stored in the order APEX, LLSB, LUSB, and RUSB. ECG and
PCG are synchronized within each site, but the sites were recorded
sequentially and are not temporally aligned with one another. Each
recording lasts 30 seconds. The four sites belong to one patient example
and are not treated as four independent patients.

\subsection{Signal processing and patient representation}
ECG is sampled at 500 Hz and band-pass filtered between 0.5 and 40 Hz.
PCG is sampled at 4,000 Hz and band-pass filtered between 25 and 400 Hz.
The preprocessing uses Butterworth filters with SciPy's order parameter
set to 4 and applies them forward and backward to avoid a phase shift.
Filtering is performed independently for each modality and recording site.

Each filtered recording is then standardized by subtracting its own mean
and dividing by its own standard deviation, with a minimum denominator of
$10^{-8}$ to avoid division by zero. This step uses no statistics from
other patients. The shared preprocessing preserves the full recording;
it does not crop, segment, pad, or resample the signals.

For each patient, the four sites are stacked into an ECG tensor of shape
$4\times15{,}000$ and a PCG tensor of shape $4\times120{,}000$, stored as
32-bit floating-point values. The first dimension represents recording
sites and the second represents time samples.

\subsection{Data checks and patient separation}
The preprocessing checks that each patient has all four required sites,
with no duplicate sites and consistent labels, LVEF values, and fold
assignments. It also checks that the source files exist and that each
signal is mono, contains finite values, and has the expected sampling
rate and number of samples. Patient identifiers must not overlap between
the development and community test cohorts.

The processed metadata contains 500 development patients and 120 community
test patients, with no overlapping identifiers. Each of the five development
folds contains 100 patients, including 12 with HFrEF. All four sites remain
together when patients are split. Model-specific training crops and
augmentation are described with the corresponding model, separately from
this shared preprocessing.

% \draftnote{Verify these details against the final report's chosen code
% revision. Add a compact cohort table, dataset citation, and a brief account
% of the checks for missing sites, inconsistent labels, and patient overlap.}
% SUGGESTION (optional): show a short ECG interval before and after filtering,
% with labeled time and amplitude axes, to illustrate the preprocessing.
% ------------- evaluation -------------
\section{Experimental Setup}
\label{sec:evaluation}

\subsection{Data partitions and model selection}
The 500-patient development cohort is divided into five predefined
patient-level folds, each containing 100 patients, including 12 with HFrEF
(Table~\ref{tab:cohorts}). The main ECG, PCG, and fusion comparison uses
five-fold out-of-fold (OOF) evaluation: in each run, four folds (400 patients)
are used for training and the remaining fold (100 patients) is held out for
validation. The procedure is rotated across all five folds so that every
development patient receives exactly one prediction from a model that was not
trained on that patient. All four recording sites from one patient remain in
the same partition, preventing site-level leakage across training and
validation~\cite{hfrefdataset}.

This five-fold protocol is distinct from the historical ECG fold-0 baseline
described in Section~\ref{sec:ecg}, which uses fold 0 as its single validation
set, and from exploratory ECG experiments that use separate training, tuning,
and scoring folds. Those results are reported with their own protocols rather
than being treated as equivalent estimates.

For the five-fold comparison, model checkpoints are selected using validation
AUROC within each fold. Decision thresholds are selected from the pooled
development OOF predictions after all five folds have been completed. Two
operating rules are saved: the threshold that maximizes development F1 and a
threshold targeting at least 90\% development specificity. For the final
community evaluation, thresholds are transferred from development without
retuning on community labels. The selected fusion architecture is also chosen
from development OOF results before its final community evaluation.

\subsection{Training and reproducibility}

The primary five-fold ECG and PCG runs use seed 42, batch size 16, an initial
learning rate of $10^{-3}$, at most 50 epochs, and early stopping with patience
10. Fold-specific seeds are derived deterministically from the base seed and
fold index. Training uses focal loss with $\gamma=2$, with the positive-class
weight set from the class distribution in the corresponding training fold;
focal loss is used to reduce the contribution of well-classified examples in
an imbalanced setting~\cite{lin2017focal}. Optimization uses AdamW
\cite{loshchilov2019adamw} with weight decay $10^{-4}$ for ECG and
$10^{-3}$ for PCG. A ReduceLROnPlateau scheduler halves the learning rate after
three epochs without improvement in validation AUROC, and gradients are clipped
to a maximum norm of 1.

The ECG five-fold runs use the complete 30-second input during both training
and validation. The PCG five-fold runs use random 15-second crops during
training while retaining the complete 30-second recording for validation. The
fusion runs use the same batch size, learning rate, maximum epoch count, and
early-stopping patience. Their pretrained ECG and PCG encoders are frozen by
default, while the fusion-specific layers are optimized with focal loss and
AdamW. Each run saves its fold assignment, random seed, model configuration,
best epoch, validation AUROC, and timing information. The comparison in
Section~\ref{sec:comparison} refers to the saved seed-42 five-fold outputs from
repository revision \texttt{d76164b}.

\subsection{Evaluation measures}

Discrimination is summarized with the area under the receiver operating
characteristic curve (AUROC), which evaluates ranking across all possible
decision thresholds~\cite{fawcett2006roc}. Because HFrEF is the minority class
in both cohorts, we also report precision--recall performance. In the saved
analysis files, the quantity labelled AUPRC is computed using average precision
(AP), i.e., the weighted summary of precision as recall increases. Precision--
recall measures are particularly informative for imbalanced binary
classification~\cite{saito2015pr}.

At a fixed threshold, sensitivity (recall) is
$TP/(TP+FN)$, specificity is $TN/(TN+FP)$, precision is $TP/(TP+FP)$, and
F1 is the harmonic mean of precision and recall. Balanced accuracy is the
average of sensitivity and specificity. Threshold-independent AUROC and AP do
not change when the operating threshold changes.

Development results are reported in two complementary ways. Fold-wise values
summarize performance separately within each held-out fold and may be averaged
across folds, while pooled OOF values are calculated once after combining the
500 patient-level OOF predictions. These quantities need not be numerically
identical because each fold is produced by a separately trained model. For
community testing, each five-fold ensemble first averages the five model
probabilities for a patient and then applies the development-derived threshold.
% ------------- ecg -------------
% Owner: Author One (ECG). Shared dataset/protocol details live elsewhere.
\section{ECG Model and Experiments}
\label{sec:ecg}

\subsection{Baseline architecture}
The ECG baseline receives all four site recordings as a single patient
input, $X_{\mathrm{ECG}}\in\mathbb{R}^{4\times15{,}000}$. A one-dimensional
convolution with 32 output channels, kernel size 15, and stride 2 is followed
by batch normalization, a rectified linear unit (ReLU), and max pooling.
Four residual blocks use 32, 64, 128, and 128 output channels. Each block
contains two convolutions with kernel size 7 and a residual connection;
the final three blocks downsample using stride 2 in the first convolution.
Batch normalization follows each convolution within a block, and channel
dropout with probability 0.1 is applied between the two convolutions.

Global average and maximum pooling produce two 128-dimensional vectors,
which are concatenated. A fully connected layer maps the resulting
256-dimensional vector to a 128-dimensional embedding, followed by ReLU
and dropout with probability 0.3. A final linear layer produces one logit
per patient. A sigmoid converts the logit to a classification score.

\subsection{Baseline training}
The historical baseline uses fold 0 for validation (100 patients, including
12 with HFrEF) and the remaining four folds for training (400 patients,
including 48 with HFrEF). It is trained on full ECG recordings using seed
42 and a batch size of 16. AdamW is used with an initial learning rate of
$10^{-3}$ and weight decay of $10^{-4}$.

The training objective is focal loss with $\alpha=0.88$ and $\gamma=2$.
In this implementation, positive examples receive a class weight of 0.88
and negative examples a weight of 0.12. The focal factor reduces the
contribution of confidently classified examples, giving relatively more
weight to difficult examples. These settings address the imbalance
between the two classes during training.

Validation AUROC controls both learning-rate reduction and checkpoint
selection. The learning rate is halved after a plateau with patience 3,
and early stopping uses patience 10 within a maximum of 50 epochs. The
saved checkpoint is the one with the highest validation AUROC, at epoch
23. Its decision threshold, approximately 0.5117, is selected by maximizing
validation F1 and then kept fixed for community testing. These settings
describe the historical baseline; the later five-fold runs are reported
separately in Section~\ref{sec:comparison}.

% \draftnote{Explain the loss and optimizer using the original ECG notebook.
% Keep these historical settings distinct from later experiments. Refer to
% Section~\ref{sec:evaluation} for the partition and threshold rules.}

\subsection{Improvement experiments}
We tested changes to the input representation, recording length, and
training procedure, and checked factors that might affect the results.

\begin{itemize}
\item \textbf{Separate-site processing and feature concatenation.}
Because the four sites were recorded at different times, we tried processing
each site with the same single-channel encoder, then concatenating the four
feature vectors before classification. This combines learned features,
rather than joining the raw recordings into one long signal. It improved
some development results, but did not outperform the historical model on
the community test set.

\item \textbf{Five-second segments.}
We trained a model on six non-overlapping five-second windows per patient,
filtering and normalizing each window separately. We tested individual
windows, averaging two or six window scores, and majority voting.
Averaging all six windows gave a validation AUROC of 0.9252, compared with
0.9564 for the historical full-recording model. On the community set, it
detected four of five HFrEF patients with four false positives; the
historical model detected all five with the same number of false positives.

\item \textbf{Fixed and random 15-second training crops.}
We compared full recordings with the middle 15 seconds and randomly chosen
15-second crops, while keeping full recordings for evaluation. Random
cropping reduced mean average precision from 0.524 to 0.476. The fixed
crop reached 0.553, but sensitivity fell from 66.7\% to 45.0\%.
Neither crop replaced the full input. These are exploratory results from
before the training-check correction described below.

\item \textbf{Heart-rate, recording-quality, and probability checks.}
Simple models using heart rate, recording quality, or both reached mean
average precisions of 0.206, 0.383, and 0.404, respectively. We also checked
whether predicted probabilities matched observed outcomes and compared
performance across signal-quality groups. Recording conditions may carry
information about the label, but these checks do not prove that the ECG
model relies on noise. No clear quality group explained the separate-site
model's improvement.

\item \textbf{Repeated architecture and threshold comparisons.}
We compared the four-channel and separate-site models over five scoring
folds with seeds 42 and 2026, using separate patients for training, tuning,
and scoring. We repeated training after fixing a preliminary model check
that had changed BatchNorm statistics and the random sequence. Alongside
the best-F1 threshold, we tested thresholds targeting at least 90\%
specificity on tuning patients. Under this second rule, separate-site
processing improved mean scoring sensitivity with both seeds, although
the size of the gain remained uncertain.

\item \textbf{Five-fold baseline evaluation and loss comparison.}
A later experiment trained the baseline architecture on four folds and
validated on the remaining fold, then averaged the five models' community
predictions. Its fold-0 result was lower than the historical result.
A controlled rerun traced this difference to two mathematically equivalent
focal-loss implementations with different numerical operations. The original
implementation exactly reproduced the historical model, including its
validation AUROC of 0.9564; the shared implementation reproduced the newer
fold-0 AUROC of 0.9110. This explains that run's difference, but does not
show that one implementation is always better.
\end{itemize}

Earlier single-channel trials also tested segmentation, weighted binary
cross-entropy, and focal loss. Their split could place recordings from the
same patient in both training and validation, so their scores are excluded
from the final comparison. Synthetic noise and additional class balancing
were considered but deferred, rather than completed experiments.

% Sources: analysis_outputs/ecg_experiments/README.md and accompanying CSVs.
% Five-second study: ea0014a:notebooks/3_ecg_5second_segments.ipynb.
% Later loss study: origin/main at d76164b, analysis_outputs/ecg_loss_comparison/seed42/.
% \draftnote{Summarize the separate-site encoder, fixed/random cropping,
% recording-quality audit, and controlled loss comparison. Use a small table
% with columns for motivation, change, result, and decision. Select the most
% informative findings instead of reproducing every training log.}

\subsection{Results and interpretation}
We retained the historical fold-0 four-channel ECG model as the project
baseline. On its 100 validation patients, it achieved an AUROC of 0.9564
and average precision of 0.7686. These are single-fold results, distinct
from the five-fold averages in Section~\ref{sec:comparison}.

In the community comparison using thresholds chosen to target at least
90\% tuning specificity, both the historical and selected separate-site
models detected all five HFrEF patients. However, the historical model
produced four false positives, compared with 14 for the separate-site
model, and had a higher AUROC (0.9983 versus 0.9913). The development gains
therefore did not lead to a better community result. This supported keeping
the historical baseline, while recognizing that the test set contains
only five positive patients.

% \draftnote{Insert verified ECG results with the exact run, split, and metric
% aggregation identified. Explain the difference between the historical
% single model and five-fold models. State what the experiments support and
% what remains uncertain; leave cross-modality conclusions to
% Section~\ref{sec:comparison}.}

% Example after a figure is available under report/figures/:
% \begin{figure}[htbp]
%   \centering
%   \includegraphics[width=0.85\linewidth]{ecg_training.pdf}
%   \caption{Training history for the specified ECG baseline run.}
%   \label{fig:ecg-training}
% \end{figure}
% ------------- pcg -------------
% Owner: Author Two (PCG).
\section{PCG Model and Experiments}
\label{sec:pcg}

\subsection{Baseline architecture}
The PCG encoder is a 1D residual CNN operating on the four-channel auscultation
signal (APEX, LLSB, LUSB, RUSB), each 30 seconds long. Because PCG is sampled
at 4{,}000~Hz against ECG's 500~Hz, the same 30-second recording is $8\times$
longer in raw samples (120{,}000 vs.\ 15{,}000), so the input tensor is
$(4, 120000)$ per patient, following the bandpass filtering and per-channel
Z-score normalization described in Section~\ref{sec:data}.

To keep the PCG and ECG branches architecturally compatible for later fusion,
the PCG stem downsamples far more aggressively than the ECG stem: a stride-4
convolution followed by a stride-4 max pool (a total of $/16$), compared to
$/4$ for ECG. This is a deliberate design choice --- after four residual
blocks with channel widths $32{\to}64{\to}128{\to}128$ (each block: two
kernel-7 convolutions, a skip connection, and 0.1 dropout between
convolutions), the resulting sequence length is approximately 469, matching
the ECG encoder's output length exactly. This alignment is what lets the two
encoders' pooled features be concatenated (or attended over) in the fusion
stage without a length-adaptation layer.

The classifier head combines global average- and max-pooled features (256-d)
into a 128-dimension embedding via a fully connected layer, ReLU, and
dropout, followed by a single linear output passed through a sigmoid for the
HFrEF probability.

\subsection{Training and improvement experiments}
\textbf{Baseline training crop and regularization.} The first training run
overfit clearly: training loss fell to near 0.0009 while validation loss rose
steadily from $\sim$0.07 to $\sim$0.24, and validation AUROC oscillated
between 0.5 and 0.65 with no upward trend. Two changes were tried in
sequence. First, increased regularization --- dropout raised from 0.3 to 0.5
and weight decay from $10^{-4}$ to $10^{-3}$ --- gave a negligible
improvement (best AUROC 0.6307 vs.\ 0.6402), showing the problem was not
simply under-regularized weights. Second, random-crop augmentation: instead
of training on the fixed 30-second recording every epoch, each training
batch uses a random 15-second crop with a new start position each time it is
sampled; validation and test still use the full 30-second recording. Because
the encoder uses adaptive pooling, no architecture change was required. This
was the fix that worked: training and validation loss converged together
around $\sim$0.035, and validation AUROC became a stable upward trend from
0.52 to 0.67 across epochs. With only $\sim$500 development patients
($\sim$65 HFrEF-positive), the crop augmentation is understood to prevent the
model from memorizing patient-specific noise in a fixed recording.

\textbf{Smaller model with augmentation, tested separately.} A further
experiment swapped in a smaller encoder (all channel widths halved relative
to the baseline) and combined several changes together: same-class mixup,
mild Gaussian noise, amplitude scaling, the full 30-second recording (no
crop), stronger weight decay, and 0.5 dropout. Because model size,
augmentation type, and crop policy all changed in the same run, the results
cannot isolate which change is responsible for the outcome --- a limitation
of the experiment design made explicit here. The result underperformed the
baseline on every metric (Section~\ref{sec:pcg-results}), so the smaller
encoder was kept only as the default for this one experiment, not adopted as
a replacement baseline.

\subsection{Results and interpretation}
\label{sec:pcg-results}
\textbf{Verification status of the numbers.} The original baseline checkpoint
is not present in the project repository --- large model files are excluded
from version control --- so its reported test metrics (AUROC 0.8748, AUPRC
0.1623, recall 1.0, precision 0.104) are \emph{historical notebook outputs},
carried forward as a fixed reference rather than reproduced by re-running a
saved model. By contrast, the five-fold cross-validation results below
\emph{are} replayed predictions from checkpoints that do exist, reported at
full precision from the saved result tables.

\textbf{Development (five-fold cross-validation).} Across all five folds (400
train / 100 validation patients each, 12 HFrEF-positive per validation fold),
per-fold validation AUROC ranged from 0.686 to 0.884. Pooled out-of-fold
predictions over all 500 development patients gave AUROC 0.757 and AUPRC
0.325 --- a more stable estimate than any single split, since every patient
is scored by a model that never trained on them.

\textbf{Community test set.} The five-model ensemble scored AUROC 0.647 /
AUPRC 0.269 on the 120-patient community test set (5 HFrEF-positive), while
the original single model reportedly reached AUROC 0.875 / AUPRC 0.162. The
two results trade off differently: the original model caught all 5 HFrEF
patients but produced 43 false positives among 115 negatives (specificity
0.626); the five-fold ensemble caught only 1--2 of 5 HFrEF patients but with
far fewer false alarms (9--12), i.e., higher specificity (0.90--0.92) at the
cost of recall. Given this trade-off and the unresolved checkpoint-
verification gap, the original model was kept as the standalone PCG
baseline, while the five-fold run was preserved as a reference experiment.

\textbf{Discrimination and the small positive class.} PCG alone is a
substantially weaker predictor of HFrEF than ECG alone (development AUROC
0.757 for PCG vs.\ 0.9564 reported for the ECG baseline). All PCG results are
also affected by the small number of positive cases: only 12 HFrEF patients
per 100-patient validation fold and just 5 in the full 120-patient community
test set. At this scale, a single changed prediction shifts test recall by
20 percentage points, so point estimates of recall, precision, and balanced
accuracy on the test set should be read as high-variance rather than
precise measurements; F1 is treated as the more informative single number
for this reason. The augmentation experiment reinforces the same pattern: it
further reduced false positives (20 vs.\ 43 on the community test) but at a
steeper recall cost (1 of 5 HFrEF patients found, AUROC 0.68, AUPRC 0.08),
and --- because several training changes were bundled together --- cannot be
used to attribute the degradation to any one factor.
% ------------- fusion -------------
% Owner: Author Three (fusion).
\section{Multimodal Fusion}
\label{sec:fusion}

\subsection{Concatenation fusion}

For each held-out fold, the fusion model loads the corresponding pretrained
ECG and PCG encoders from the five-fold seed-42 runs. The encoder associated
with a given fold was trained without that fold, so the same held-out patients
remain reserved for validation during fusion training. The ECG and PCG
architectures are described in Sections~\ref{sec:ecg} and~\ref{sec:pcg}.

Each encoder produces a 128-dimensional patient embedding. In the
concatenation model, these two embeddings are placed side by side to form a
256-dimensional feature vector. This is feature-level fusion: the raw ECG and
PCG signals are not added, averaged, or directly mixed. The combined vector is
passed through a small multilayer perceptron consisting of a linear layer from
256 to 128 units, ReLU activation, dropout with probability 0.3, and a final
linear layer that produces one classification logit.

The pretrained ECG and PCG encoders are frozen during this experiment and are
kept in evaluation mode. Therefore, only the fusion classifier is optimized.
Although concatenation does not explicitly model cross-modal attention, the
classifier can still learn weighted combinations of ECG and PCG features after
the two representations have been joined.

\subsection{Attention fusion}
The attention model starts from the same fold-specific pretrained ECG and PCG
encoders and keeps them frozen. Their 128-dimensional embeddings are projected
separately into a common 128-dimensional attention space. The projected ECG
embedding is treated as one modality token and the projected PCG embedding as
a second modality token.

The two tokens are processed by multi-head self-attention with four heads.
This allows the ECG token to attend to both ECG and PCG information, and
likewise allows the PCG token to attend to both modalities. The attention
output is combined with the original tokens through a residual connection and
layer normalization. A feed-forward block is then applied, followed by a
second residual connection and layer normalization.

The two updated 128-dimensional tokens are flattened to a 256-dimensional
vector and passed to the final multilayer perceptron classifier. This mechanism
models interaction between whole-patient ECG and PCG embeddings. It should not
be interpreted as explicit beat-level or sample-level temporal alignment
between the two raw signals.

\subsection{Development selection and community evaluation}
Both fusion architectures were evaluated with the same five-fold
out-of-fold development protocol. In each fold, the fusion head was trained on
four development folds and validated on the remaining fold, using the
corresponding pretrained ECG and PCG encoders. The five validation partitions
were then combined so that every development patient received an out-of-fold
prediction from a model that had not trained on that patient.

On pooled development predictions, concatenation achieved an AUROC of 0.9034
and average precision of 0.5735, compared with 0.8357 and 0.4962 for the
attention model. Based on this development comparison, concatenation was
selected as the final fusion architecture. Its maximum-development-F1
threshold was 0.5773.

For the community cohort, each of the five concatenation fold models produced
a probability for every patient. The five probabilities were averaged to form
the ensemble prediction, and the development-selected threshold was applied
without retuning on the community data. This ensemble achieved an AUROC of
0.9948 and average precision of 0.9250. At the 0.5773 threshold it detected
four of five HFrEF patients, with one false negative, one false positive, and
114 true negatives. Attention fusion was not used for the final community
evaluation because the fusion architecture had already been selected from the
development out-of-fold results.

% ------------- comparison discussion -------------
% Shared: each owner verifies their rows; one editor integrates the text.
\section{Comparative Results and Discussion}
\label{sec:comparison}

\subsection{Comparison of modalities and fusion strategies}

Table~\ref{tab:comparison} compares the published five-fold results from
Ma et al.~\cite{ma2026hierarchical} with our saved seed-42 development
predictions. Our entries are the mean and sample standard deviation across the
five held-out folds, each containing 100 patients, including 12 with HFrEF.
The published row is reproduced as reported rather than obtained by retraining
LGC$^{2}$-Net; implementation details for the published protocol are also
available in the authors' official repository~\cite{lgc2implementation}.
Accordingly, the two sets of results should be treated as contextual rather
than as a controlled head-to-head replication.

For threshold-dependent development metrics, each of our models uses a
threshold selected from the pooled OOF predictions by maximizing development
F1. Those thresholds are then applied to the individual held-out folds when the
fold-wise sensitivity, specificity, balanced accuracy, and overall accuracy are
summarized. Because the same pooled OOF predictions are used to choose these
development operating points, the threshold-dependent development values are
descriptive model-selection results rather than estimates from an untouched
test cohort.

% Results source: origin/main, commit d76164b1ada6b81baececf98fc4fbd12f4dd9552.
% Directory: analysis_outputs/five_fold_cross_validation/.
% Runs: ecg_seed42, pcg_seed42, fusion_concat_seed42, fusion_attention_seed42.
% Recomputed from development_oof_predictions.csv, grouped by Fold.
% Thresholds: development_summary.csv, maximum_development_f1.
% SD uses ddof=1. AUROC is computed within each fold before averaging.
% These are the five-fold runs, not the historical fold-0 baseline models.

\begin{table}[htbp]
\centering
\caption{Published five-fold results and our development validation results.
Our entries are mean $\pm$ sample standard deviation across folds.
All metrics except AUROC are percentages. Protocols differ as described in the
text.}
\label{tab:comparison}
\begingroup
\setlength{\tabcolsep}{4pt}
\resizebox{\textwidth}{!}{%
\begin{tabular}{lccccc}
\toprule
Model & Sensitivity & Specificity & Balanced acc. & Overall acc. & AUROC \\
\midrule
LGC$^{2}$-Net~\cite{ma2026hierarchical}
& $90.00\pm6.97$ & $92.73\pm2.62$ & $91.36\pm4.12$ & $92.40\pm2.70$ & $0.960\pm0.020$ \\
\midrule
Our ECG
& $55.00\pm29.81$ & $92.73\pm8.75$ & $73.86\pm11.21$ & $88.20\pm4.82$ & $0.9121\pm0.0230$ \\
Our PCG
& $38.33\pm32.60$ & $91.36\pm9.73$ & $64.85\pm12.60$ & $85.00\pm5.79$ & $0.7767\pm0.0853$ \\
Our concatenation
& $63.33\pm20.92$ & $94.55\pm2.94$ & $78.94\pm9.34$ & $90.80\pm1.64$ & $0.9288\pm0.0311$ \\
Our attention
& $65.00\pm29.11$ & $92.50\pm2.85$ & $78.75\pm13.24$ & $89.20\pm1.48$ & $0.9100\pm0.0704$ \\
\bottomrule
\end{tabular}
}
\endgroup
\end{table}

Concatenation achieved the highest mean fold AUROC among our saved runs
(0.9288), followed by ECG (0.9121) and attention (0.9100). Attention showed a
slightly higher mean sensitivity than concatenation but lower specificity and
a larger fold-to-fold AUROC standard deviation. These differences describe the
observed runs; no statistical test in this project establishes that the AUROC
differences between our models are significant.

Mean fold metrics should be distinguished from metrics calculated after pooling
all 500 OOF predictions. In the pooled fusion comparison, concatenation
achieved AUROC 0.9034 and average precision 0.5735, while attention achieved
0.8357 and 0.4962. Mean fold average precision for concatenation and attention
was much closer (0.6420 and 0.6429). This difference is plausible because each
held-out fold is produced by a separately trained model whose probability scale
can differ from the others; pooled OOF metrics therefore answer a different
aggregation question than the mean of fold-wise metrics. AUROC and
precision--recall measures should also be interpreted together in this
imbalanced setting~\cite{fawcett2006roc,saito2015pr}.

\subsection{Community-test results}

Table~\ref{tab:community-comparison} reports community results separately from
development results. For our five-fold ensembles, each fold model predicts all
120 community patients, the five probabilities for each patient are averaged,
and the development-derived operating threshold is then applied. This mirrors
the saved ensemble procedure rather than the published LGC$^{2}$-Net
aggregation protocol~\cite{lgc2implementation}.

\begin{table}[htbp]
\centering
\caption{Community-cohort results. Sensitivity, specificity, and balanced
accuracy are percentages. Published values and our ensemble results use
different model aggregation procedures.}
\label{tab:community-comparison}
\begingroup
\setlength{\tabcolsep}{4pt}
\resizebox{\textwidth}{!}{%
\begin{tabular}{lcccc}
\toprule
Model & Sensitivity & Specificity & Balanced acc. & AUROC \\
\midrule
LGC$^{2}$-Net (reported)~\cite{ma2026hierarchical} & 100.00 & 92.52 & 96.26 & 0.990 \\
\midrule
Our ECG ensemble & 60.00 & 100.00 & 80.00 & 0.9896 \\
Our PCG ensemble & 20.00 & 92.17 & 56.09 & 0.6470 \\
Our concatenation ensemble & 80.00 & 99.13 & 89.57 & 0.9948 \\
\midrule
Our historical ECG single model & 100.00 & 96.52 & 98.26 & 0.9983 \\
\bottomrule
\end{tabular}
}
\endgroup
\end{table}

The concatenation ensemble detected four of the five HFrEF patients and
produced one false positive among 115 non-HFrEF patients. Its average precision
was 0.9250. The historical ECG single model detected all five positive patients
but produced four false positives. These observations show a sensitivity--
specificity trade-off rather than proving that fusion is universally superior
to ECG alone. The community cohort contains only five HFrEF cases, so one
additional missed or detected case changes sensitivity by 20 percentage
points. The small positive class also makes precision--recall estimates
particularly sensitive to the ranking of individual positive
patients~\cite{saito2015pr}.

Attention fusion was not included in the final community evaluation because
concatenation had already been selected from development OOF results. Running
multiple alternative fusion architectures on the community cohort and then
selecting the best one would use the community labels for model selection,
weakening its role as a final evaluation set.

\subsection{Limitations}

Several limitations constrain the interpretation of the results. First, our
results are not directly comparable with the published LGC$^{2}$-Net values
because the architectures, training procedures, threshold rules, and model
aggregation procedures differ~\cite{ma2026hierarchical,lgc2implementation}.
Second, development thresholds are selected from pooled OOF predictions, so
threshold-dependent development metrics are part of model selection rather
than independent test estimates. Third, the community cohort contains only
five HFrEF patients, creating substantial uncertainty in sensitivity,
precision, and related metrics. Fourth, the same community cohort had already
been inspected during earlier ECG experiments, so the final community
evaluation should not be described as a completely untouched external
validation.

The fusion implementation does preserve patient-level fold alignment: each
fusion fold reuses the ECG and PCG checkpoints from the corresponding
five-fold baseline run, and those encoders were trained on the same four
training folds while the remaining fold stayed held out. The encoders are also
frozen during the reported fusion training. This reduces leakage within the
development comparison, but it does not remove the broader limitations above.
The results therefore support concatenation as the strongest fusion approach
among the architectures tested in this project, while larger and genuinely
unseen external cohorts would be required before making stronger claims about
generalization or clinical utility.

% SUGGESTION (before submission): match each fusion run to its exact pretrained
% ECG and PCG checkpoints and reconcile the notebook and prediction-file results.
% ------------- conclusion -------------
\section{Conclusion}
\label{sec:conclusion}

This project investigated whether combining ECG and PCG could improve HFrEF
detection compared with using either modality alone. Across the development
cohort, ECG was the stronger unimodal signal, while simple ECG--PCG feature
concatenation provided the strongest fusion result. On pooled out-of-fold
predictions, concatenation achieved an AUROC of 0.9034 and average precision of
0.5735, compared with 0.8661 and the corresponding lower precision--recall
performance for ECG alone. The attention-based model did not improve on simple
concatenation, indicating that additional fusion complexity was not beneficial
for this dataset and training setup.

The selected concatenation ensemble also performed strongly on the community
cohort, detecting four of the five HFrEF patients with one false positive.
However, the very small number of positive community cases, together with the
fact that this cohort had been inspected during earlier experiments, limits how
strongly these results can be generalized. The findings therefore support ECG
and PCG as potentially complementary modalities, but they do not establish that
fusion is universally superior to a strong ECG model.

Future work should evaluate the selected fusion approach on a larger,
genuinely unseen external cohort and examine whether more targeted multimodal
strategies, such as limited encoder fine-tuning or finer-grained temporal
interaction, can improve attention-based fusion without unnecessary model
complexity. Explainability and efficiency analyses would also help clarify how
each modality contributes to the final prediction and whether the approach is
practical for deployment.



% ------------- References -------------
% References are written here directly; no separate bibliography file is needed.
\begin{thebibliography}{9}

\bibitem{hfrefdataset}
F. Ma et al.,
\textit{Multi-Channel ECG and PCG Dataset for Heart Failure with Reduced
Ejection Fraction (HFrEF) Detection}, Zenodo.
\url{https://zenodo.org/records/16934966}.

\bibitem{ma2026hierarchical}
F. Ma et al.,
``Hierarchical fusion of electrocardiogram and phonocardiogram data
improves heart failure detection,''
\textit{Patterns}, vol. 7, no. 3, article 101448, 2026.
\url{https://doi.org/10.1016/j.patter.2025.101448}.

\bibitem{lgc2implementation}
F. Ma et al., LGC$^{2}$-Net official implementation,
five-fold and community evaluation scripts. Accessed September 24, 2026.
\url{https://github.com/Ezio660914/LGC2_Net_HFrEF}.

\bibitem{lin2017focal}
T.-Y. Lin, P. Goyal, R. Girshick, K. He, and P. Dollár,
``Focal Loss for Dense Object Detection,''
in \textit{Proceedings of the IEEE International Conference on Computer Vision (ICCV)},
pp. 2980--2988, 2017.
\url{https://doi.org/10.1109/ICCV.2017.324}.

\bibitem{loshchilov2019adamw}
I. Loshchilov and F. Hutter,
``Decoupled Weight Decay Regularization,''
in \textit{7th International Conference on Learning Representations (ICLR)}, 2019.
\url{https://openreview.net/forum?id=Bkg6RiCqY7}.

\bibitem{fawcett2006roc}
T. Fawcett,
``An Introduction to ROC Analysis,''
\textit{Pattern Recognition Letters}, vol. 27, no. 8, pp. 861--874, 2006.
\url{https://doi.org/10.1016/j.patrec.2005.10.010}.

\bibitem{saito2015pr}
T. Saito and M. Rehmsmeier,
``The Precision--Recall Plot Is More Informative than the ROC Plot When Evaluating Binary Classifiers on Imbalanced Datasets,''
\textit{PLOS ONE}, vol. 10, no. 3, e0118432, 2015.
\url{https://doi.org/10.1371/journal.pone.0118432}.

\bibitem{vaswani2017attention}
A. Vaswani et al.,
``Attention Is All You Need,''
in \textit{Advances in Neural Information Processing Systems 30}, 2017.
\url{https://arxiv.org/abs/1706.03762}.


\end{thebibliography}
\end{document}
